\section{Results and Discussion}
\label{sec:results}

This section reports the revised empirical evidence for the communication-aware adaptive weighting design. The core component values come from one controlled 50-trial experiment with deterministic seeds $2$--$51$. The Cardinality-critical mode is evaluated in a separate deterministic 50-trial run using the same seeds, scenario generator, communication protocol, and no-stabilization configuration. This design preserves matched scenario realizations while avoiding reuse of the older stabilized FID-FIA result.

\subsection{Main result in the tiered-drop GA scenario}

Table~\ref{tab:main-consensus-stats} compares the two no-stabilization operating modes with the most informative lower-complexity steps on the same deterministic trial set. Covariance weighting improves over Fixed Metropolis, but realized link quality produces the dominant reduction under heterogeneous packet loss. Adding covariance to the link-only rule provides a further cardinality-dispersion reduction, while Balanced gives the lowest OSPA consensus error and matched localization disagreement. Cardinality-critical adds FID-FIA only to the existence branch and gives the lowest cardinality dispersion.

\begin{table}[t]
  \centering
  \caption{Main network-disagreement comparison under tiered heterogeneous packet loss. Values are mean $\pm$ standard deviation over 50 deterministic trials with seeds $2$--$51$. Both operating modes use instantaneous adaptive weights without EMA smoothing or final-weight floors. Lower values are better.}
  \label{tab:main-consensus-stats}
  \begin{tabular}{lccc}
    \toprule
    Arm & OSPA err. & Loc. disag. & Card. disp. \\
    \midrule
    Fixed Metropolis & $2.383 \pm 0.176$ & $2.263 \pm 0.340$ & $0.650 \pm 0.223$ \\
    Covariance only & $1.983 \pm 0.120$ & $1.877 \pm 0.240$ & $0.372 \pm 0.114$ \\
    Link quality only & $1.712 \pm 0.049$ & $1.474 \pm 0.053$ & $0.102 \pm 0.026$ \\
    Covariance + link & $1.708 \pm 0.046$ & $1.483 \pm 0.059$ & $0.092 \pm 0.025$ \\
    Balanced mode & $\mathbf{1.696} \pm 0.046$ & $\mathbf{1.461} \pm 0.053$ & $0.095 \pm 0.025$ \\
    Cardinality-critical mode & $1.713 \pm 0.049$ & $1.590 \pm 0.167$ & $\mathbf{0.062} \pm 0.016$ \\
    \bottomrule
  \end{tabular}
\end{table}

Relative to Fixed Metropolis, the Balanced mode reduces OSPA consensus error by $28.8\%$, matched localization disagreement by $35.5\%$, and cardinality dispersion by $85.4\%$. The more important diagnostic is incremental: the combined branch-aware refinement improves the shared three-factor arm by $0.0145$ in OSPA and $0.0202$ in localization disagreement, with improvements in 50/50 and 49/50 paired trials, respectively. Cardinality-critical then reduces cardinality dispersion by a further $34.8\%$ relative to Balanced, with lower dispersion in 45/50 matched seeds. This gain is specialized rather than uniform: OSPA consensus error increases by $1.0\%$ and localization disagreement by $8.9\%$.

Table~\ref{tab:main-local} checks whether stronger agreement is obtained by driving all nodes toward an inaccurate common estimate. Balanced remains close to the covariance-link backbone on all three local metrics. Cardinality-critical improves local E-OSPA and reduces local cardinality error by $26.2\%$ relative to Balanced, but increases local RMSE by $6.6\%$. This mirrors the network-level spatial-versus-cardinality tradeoff.

\begin{table}[t]
  \centering
  \caption{Truth-referenced local tracking metrics for the same 50 deterministic seeds as Table~\ref{tab:main-consensus-stats}. Each trial is averaged over sensors before the across-trial mean and standard deviation are computed.}
  \label{tab:main-local}
  \begin{tabular}{lccc}
    \toprule
    Arm & E-OSPA & RMSE & CardErr \\
    \midrule
    Fixed Metropolis & $2.781 \pm 0.129$ & $1.630 \pm 0.049$ & $1.364 \pm 0.260$ \\
    Covariance only & $2.556 \pm 0.106$ & $\mathbf{1.567} \pm 0.049$ & $0.981 \pm 0.155$ \\
    Link quality only & $2.074 \pm 0.053$ & $1.626 \pm 0.045$ & $0.300 \pm 0.048$ \\
    Covariance + link & $2.064 \pm 0.051$ & $1.633 \pm 0.044$ & $0.280 \pm 0.045$ \\
    Balanced mode & $2.072 \pm 0.051$ & $1.636 \pm 0.045$ & $0.283 \pm 0.046$ \\
    Cardinality-critical mode & $\mathbf{2.030} \pm 0.042$ & $1.744 \pm 0.161$ & $\mathbf{0.209} \pm 0.026$ \\
    \bottomrule
  \end{tabular}
\end{table}

\subsection{Computational cost}

Table~\ref{tab:runtime-cost} reports filtering/fusion wall-clock time. Scenario generation, communication sampling, plotting, and metric evaluation are excluded. The core component arms were measured together in one environment and add approximately $8$--$12\%$ over Fixed Metropolis. Cardinality-critical was run separately under Octave; its absolute time confirms that FID-FIA dominates computation, but a cross-environment relative multiplier is not reported.

\begin{table}[t]
  \centering
  \caption{Filtering/fusion runtime over 50 deterministic trials. Relative runtime is computed against Fixed Metropolis only for arms measured in the same component run. The Cardinality-critical value is from the separate Octave validation and is not used for a cross-environment ratio.}
  \label{tab:runtime-cost}
  \begin{tabular}{lccc}
    \toprule
    Arm & Runtime (s) & Runtime/step (s) & Relative runtime \\
    \midrule
    Fixed Metropolis & $4.387 \pm 1.043$ & $0.0439$ & $1.000\times$ \\
    Covariance only & $4.751 \pm 1.124$ & $0.0475$ & $1.084\times$ \\
    Link quality only & $4.719 \pm 1.105$ & $0.0472$ & $1.078\times$ \\
    Covariance + link & $4.835 \pm 1.188$ & $0.0484$ & $1.115\times$ \\
    Balanced mode & $4.793 \pm 1.155$ & $0.0479$ & $1.099\times$ \\
    Cardinality-critical mode & $144.057 \pm 5.003$ & $1.4406$ & n/a \\
    \bottomrule
  \end{tabular}
\end{table}

\subsection{Factor-by-factor ablation}

Table~\ref{tab:ablation} separates isolated factor effects from the progressive construction of the retained method. The component matrix makes each comparison explicit: the first block tests covariance, realized link quality, and existence confidence independently, while the second block starts from the covariance-link backbone, adds existence confidence with shared weights, and then applies the branch-aware refinement used by Balanced mode. This final component comprises spatial/existence branch decoupling together with weak structure-aware modulation. The last block adds existence-only FID-FIA to obtain the Cardinality-critical mode.

\begin{table}[t]
  \centering
  \scriptsize
  \setlength{\tabcolsep}{2.2pt}
  \caption{Component ablation under tiered heterogeneous packet loss. Values are mean $\pm$ standard deviation over 50 deterministic trials with matched seeds. Cov., Link, Exist., Branch, and FID denote covariance quality, realized link quality, existence confidence, the branch-aware refinement, and existence-only FID-FIA modulation, respectively. The Branch component jointly includes spatial/existence decoupling and weak structure-aware modulation.}
  \label{tab:ablation}
  \begin{tabular*}{\linewidth}{@{\extracolsep{\fill}}lccccc@{\hspace{6pt}}ccc@{}}
    \toprule
    & \multicolumn{5}{c}{Active components} & \multicolumn{3}{c}{Network disagreement} \\
    \cmidrule(lr){2-6}\cmidrule(l){7-9}
    Variant & Cov. & Link & Exist. & Branch & FID & OSPA & Loc. & Card. \\
    \midrule
    \multicolumn{9}{c}{\textit{Reference}} \\
    \cmidrule(lr){1-9}
    Fixed Metropolis
      & -- & -- & -- & -- & --
      & $2.383 \pm 0.176$ & $2.263 \pm 0.340$ & $0.650 \pm 0.223$ \\
    \midrule
    \multicolumn{9}{c}{\textit{Single-factor diagnostics}} \\
    \cmidrule(lr){1-9}
    Covariance only
      & $\checkmark$ & -- & -- & -- & --
      & $1.983 \pm 0.120$ & $1.877 \pm 0.240$ & $0.372 \pm 0.114$ \\
    Link quality only
      & -- & $\checkmark$ & -- & -- & --
      & $1.712 \pm 0.049$ & $1.474 \pm 0.053$ & $0.102 \pm 0.026$ \\
    Existence confidence only
      & -- & -- & $\checkmark$ & -- & --
      & $2.057 \pm 0.143$ & $1.857 \pm 0.228$ & $0.447 \pm 0.152$ \\
    \midrule
    \multicolumn{9}{c}{\textit{Progressive retained design}} \\
    \cmidrule(lr){1-9}
    Covariance + link
      & $\checkmark$ & $\checkmark$ & -- & -- & --
      & $1.708 \pm 0.046$ & $1.483 \pm 0.059$ & $0.092 \pm 0.025$ \\
    + existence (shared)
      & $\checkmark$ & $\checkmark$ & $\checkmark$ & -- & --
      & $1.710 \pm 0.045$ & $1.481 \pm 0.055$ & $0.093 \pm 0.026$ \\
    Balanced mode
      & $\checkmark$ & $\checkmark$ & $\checkmark$ & $\checkmark$ & --
      & $\mathbf{1.696} \pm 0.046$ & $\mathbf{1.461} \pm 0.053$ & $0.095 \pm 0.025$ \\
    \midrule
    \multicolumn{9}{c}{\textit{Existence-branch extension}} \\
    \cmidrule(lr){1-9}
    Cardinality-critical mode
      & $\checkmark$ & $\checkmark$ & $\checkmark$ & $\checkmark$ & $\checkmark$
      & $1.713 \pm 0.049$ & $1.590 \pm 0.167$ & $\mathbf{0.062} \pm 0.016$ \\
    \bottomrule
  \end{tabular*}
\end{table}

The detailed results support five bounded conclusions. First, realized link quality is the dominant factor in this packet-loss scenario. Second, covariance provides a useful independent signal and gives a further cardinality benefit when combined with link quality. Third, existence confidence has little isolated incremental effect on the covariance-link backbone. Fourth, the combined branch-aware refinement supplies a small, repeatable spatial benefit. Fifth, existence-only FID-FIA provides a substantial additional cardinality benefit but does not improve the spatial disagreement metrics. The table intentionally does not separate branch decoupling from its weak structure-aware modulation.

\subsection{Scope of the revised evidence}

The revised Balanced and Cardinality-critical configurations are confirmed here only in the main tiered heterogeneous packet-loss scenario. Earlier ideal-communication, communication-level, and arithmetic-average probes used the retired stabilization settings and are therefore not used as evidence for the current modes. Repeating those supporting studies with the revised configuration is left as a targeted validation step rather than mixing incompatible operating points in the same claim.
